% Generated draft: requires mathematical and source review.
\documentclass[11pt]{article}
\usepackage[margin=25mm]{geometry}
\usepackage{fontspec}
\setmainfont{Latin Modern Roman}
\usepackage{amsmath,amssymb,mathtools,unicode-math}
\setmathfont{Latin Modern Math}
\usepackage{enumitem,needspace,xurl,hyperref}
\setlength{\parindent}{0pt}
\setlength{\parskip}{5pt}
\setlength{\emergencystretch}{3em}
\newcommand{\N}{\mathbb{N}}
\newcommand{\Z}{\mathbb{Z}}
\newcommand{\Q}{\mathbb{Q}}
\newcommand{\R}{\mathbb{R}}
\newcommand{\C}{\mathbb{C}}
\newcommand{\F}{\mathbb{F}}
\newcommand{\A}{\mathbb{A}}
\newcommand{\PP}{\mathbb{P}}
\newcommand{\E}{\mathbb{E}}
\newcommand{\eps}{\varepsilon}
\newcommand{\dd}{\,\mathrm{d}}
\begin{document}
\section*{UnsolvedMath ID: 20000425}
\textbf{Source problem code:} catalogue entry
\subsection{Definitions and mathematical statement}
\textbf{Original question (C. Hall).} Let $X(\ell)$ be a modular curve. What can be said about the splitting field of the $\ell$-torsion of $J(X(\ell))$? Hall's comment says this information is needed to discuss the $L$-function of $E\to X(\ell)$ modulo $\ell$, as a polynomial in $\F_\ell[T]$. The family $E$, ground field, level convention and desired information about the field are not specified. The OCR backticks are restored as $\ell$; the supplied verification identifies this as item 14, although the extraction and source-item metadata say 4.

\textbf{Attached unreviewed principal-level interpretation.} Let $\ell$ be prime, $K_\ell=\Q(\zeta_\ell)$, and let $X(\ell)$ be the standard geometrically connected principal-level modular curve over this cyclotomic field, with a fixed Weil-pairing component. Put $J_\ell=J(X(\ell))$, $g=\dim J_\ell$, and
\[\rho_\ell:G_{K_\ell}\longrightarrow\operatorname{Aut}_{\F_\ell}(J_\ell[\ell]),\qquad L_\ell=\overline K_\ell^{\ker\rho_\ell}=K_\ell(J_\ell[\ell]).\]
The aid reports that $L_\ell/K_\ell$ is unramified outside primes above $\ell$ and that the polarization, with $\mu_\ell\subset K_\ell$, places its image in $\mathrm{Sp}_{2g}(\F_\ell)$.

Let $\rho_\ell^{\mathrm{ss}}$ be its semisimplification and $L_\ell^{\mathrm{ss}}$ the field fixed by its kernel. The kernel and field are independent of the choice of a model for the semisimplification. The reported canonical tower and bounds are
\[K_\ell\subseteq L_\ell^{\mathrm{ss}}\subseteq L_\ell,\qquad\operatorname{Gal}(L_\ell/L_\ell^{\mathrm{ss}})\text{ an }\ell\text{-group},\qquad [L_\ell:L_\ell^{\mathrm{ss}}]\le\ell^{g^2},\]
\[[L_\ell:K_\ell]\mid |\mathrm{Sp}_{2g}(\F_\ell)|=\ell^{g^2}\prod_{i=1}^g(\ell^{2i}-1).\]
The $\ell^{g^2}$ bound is a Sylow bound, not a claim that it is attained by this particular torsion image.

The aid asserts the following equivalent conditions: $L_\ell/K_\ell$ is an $\ell$-extension; $J_\ell[\ell]^{\mathrm{ss}}$ is a direct sum of trivial representations; and every good local Frobenius polynomial away from $\ell$ satisfies
\[P_v(T)\equiv(1-T)^{2g}\pmod\ell.\]
It also says a density-one set of such good places suffices, with density interpreted for primes of $K_\ell$. For $\ell=2,3,5$ it reports $g=0$, hence $L_\ell=K_\ell$.

The reduction separates semisimple data from the upper extension-class layer: local characteristic polynomials alone do not recover that layer. It does not determine the full torsion image or splitting field in general, nor identify the unspecified mod-$\ell$ $L$-function in Hall's motivation. All ramification, group-image and criterion claims are retained as research-aid assertions requiring review.
\subsection{Short English statement}
Hall asks about the torsion splitting field of a modular-curve Jacobian. The aid chooses principal level over the cyclotomic field, gives a semisimple tower and symplectic degree bounds, and characterizes an $\ell$-extension through Frobenius polynomials. These tests do not recover the extension-class layer or the full field.
\subsection{Sources}
\begin{itemize}
\item Local UnsolvedMath registry, permanent ID 20000425, source code catalogue entry; \texttt{aim-algebraic-number-theory-notes.json}, notes[76], tag \texttt{problem}.
\item AIM workshop \emph{The Tate conjecture}, C. Hall's question: \url{https://aimath.org/WWN/tateconjecture/tateconjecture.pdf}. Source extraction labels it 4; supplied verification says official item 14 and restores $\ell$ from OCR backticks. No PDF was consulted here.
\item Attached machine-generated research attempt 1: result \texttt{partial\_result}, problem status at run \texttt{open}, rigor label \texttt{proved}, novelty confidence low. It is not peer reviewed and requires independent verification before citation. Original literature field empty; no authoritative publication references are supplied.
\end{itemize}
\end{document}
